% -*- latex -*-
%%%%%%%%%%%%%%%%%%%%%%%%%%%%%%%%%%%%%%%%%%%%%%%%%%%%%%%%%%%%%%%%
%%%%%%%%%%%%%%%%%%%%%%%%%%%%%%%%%%%%%%%%%%%%%%%%%%%%%%%%%%%%%%%%
%%%%
%%%% This text file is part of the source of 
%%%% `The Art of HPC, vol 1: The Science of Computing'
%%%% by Victor Eijkhout, copyright 2012-2023
%%%%
%%%% This book is distributed under a Creative Commons Attribution 3.0
%%%% Unported (CC BY 3.0) license and made possible by funding from
%%%% The Saylor Foundation \url{http://www.saylor.org}.
%%%%
%%%% gaussparallel.tex: LU factorization in parallel
%%%%
%%%%%%%%%%%%%%%%%%%%%%%%%%%%%%%%%%%%%%%%%%%%%%%%%%%%%%%%%%%%%%%%
%%%%%%%%%%%%%%%%%%%%%%%%%%%%%%%%%%%%%%%%%%%%%%%%%%%%%%%%%%%%%%%%

The matrix-vector and matrix-product are easy to parallelize in one sense.
The elements of the output can all be computed independently and in any order,
so we have many degrees of freedom in parallelizing the algorithm.
This does not hold for computing an LU factorization, or solving a linear system
with the factored matrix.

\Level 1 {Solving a triangular system}
\index{LU factorization!solution in parallel|(}

The solution of a triangular system $y=L\inv x$ (with $L$~is lower triangular)
is a matrix-vector operation, so
it has its $O(N^2)$ complexity in common with the matrix-vector product.
However, unlike the product operation, this solution process contains a recurrence
relation between the output elements: 
\[ y_{i} = \ell_{ii}\inv ( x_i-\sum_{j<i} \ell_{ij}x_j ). \]
This means that parallelization is not trivial.
In the case of a sparse matrix special strategies may be possible; see section~\ref{sec:parallel-implicit}.
Here we will make a few remarks about general, dense case.

Let us assume for simplicity that communication takes no time, and
that all arithmetic operations take the same unit time.
First we consider the matrix distribution by rows,
meaning that processor~$p$ stores the elements~$\ell_{p*}$.
With this we can implement the triangular solution as:
\begin{itemize}
\item Processor 1 solves $y_1=\ell_{11}\inv x_1$ and sends its value to the next processor.
\item In general, processor $p$ gets the values $y_1,\ldots,y_{p-1}$
  from processor~$p-\nobreak1$, and computes~$y_p$;
\item Each processor $p$ then sends $y_1,\ldots,y_p$ to $p+1$.
\end{itemize}

\begin{exercise}
  Show that this algorithm takes time $2N^2$, just like the sequential algorithm.
\end{exercise}

This algorithm has each processor passing all computed $y_i$ values to its successor, in a pipeline fashion.
However, this means that processor~$p$ receives~$y_1$ only at the last moment, whereas that value was computed
already in the first step. We can formulate the solution algorithm in such a way that computed elements
are made available as soon as possible:
\begin{itemize}
\item Processor~1 solve $y_1$, and sends it to all later processors.
\item In general, processor~$p$ waits for individual messages with values $y_q$ for~$q<p$.
\item Processor $p$ then computes $y_p$ and sends it to processors~$q$ with $q>p$.
\end{itemize}

Under the assumption that communication time is negligible, this algorithm can be much faster.
For instance, all processors~$p>1$ receive~$y_1$ simultaneously, and can compute $\ell_{p1}y_1$
simultaneously.

\begin{exercise}
  Show that this algorithm variant takes a time $O(N)$,
  if we ignore communication.
  What is the cost if we incorporate communication in the cost?
\end{exercise}

\begin{exercise}
  Now consider the matrix distribution by columns: processor~$p$ stores~$\ell_{*p}$.
  Outline the triangular solution algorithm with this distribution, and show that the
  parallel solve time is~$O(N)$.
\end{exercise}

\index{LU factorization!solution in parallel|)}

\Level 1 {Factorization, dense case}
\index{LU factorization!computation in parallel!dense|(}
\label{sec:LUscaling}

Analysis of Gaussian elimination, or LU factorization
is a little tricky.
First of all, 
Gaussian elimination can be performed in right-looking,
left-looking (Doolittle) and Crout variants;
see section~\ref{sec:LUvariants} and~\cite{TSoPMC}.
The choice of algorithm,
while they all compute the same result,
make have computational ramifications.

Let's limit ourselves to the 
regular right-looking Gaussian elimination
of section~\ref{sec:lu-algorithm},
reproduced in figure~\ref{eq:LU-algorithm-again}.

\begin{figure}[t]
  \begin{tabbing}
    \kern10pt\=\kern10pt\=\kern10pt\=\kill
    for\={} $k=1,n-1$:\\
    \>for\={} $i=k+1$ to $n$:\\
    \>\> $a_{ik}\leftarrow a_{ik}/a_{kk}$\\
    \>for\={} $i=k+1$ to $n$:\\
    \>\>for\={} $j=k+1$ to $n$:\\
    \>\>\>$a_{ij}\leftarrow a_{ij}-a_{ik}*a_{kj}$
  \end{tabbing}
    \label{eq:LU-algorithm-again}
    \caption{LU factorization algorithm}
\end{figure}

The number of operations in this is
\[ \sum_{k=1}^n (\sum_{i=k+1}^n 1 + \sum_{i=k+1}^n\sum_{j=k+1}^n 1 ) =n^3/3 +O(n^2). \]
The exact calculation does not concern us here;
by analogy to integrals you'll see the justice in that factor of~$1/3$
on the leading term.
Also, we will ignore lower order terms throughout.

If you do a few steps of the algorithm by hand,
you see that
\[ 
\begin{array}{l}
  a_{22} \leftarrow a_{22} - a_{21}*a_{11}\inv a_{12}\\
  \cdots\\
  a_{33} \leftarrow a_{33} - a_{32}*a_{22}\inv a_{23}\\
\end{array}
\]
meaning that there is a chain of dependent operations.

\begin{exercise}
  \label{ex:gauss-critical}
  Follow this argument through.
  Argue that there is a non-trivial \indexterm{critical path}
  in the sense of section~\ref{sec:critical-path}.
  What is its length?

  In the analysis of the critical path section,
  what does this critical path imply for the
  minimum parallel execution time and bounds on speedup?
\end{exercise}

Now let's take a closer look at how we would actually parallelize this.
The major work step is the subblock update:

\begin{tabbing}
  \kern20pt\=\kern10pt\=\kern10pt\=\kern10pt\=\kill
  for\={} $i=k+1$ to $n$:\\
  \>for\={} $j=k+1$ to $n$:\\
  \>\>$a_{ij}\leftarrow a_{ij}-a_{ik}*a_{kj}$
\end{tabbing}

In the $k=1$ step, this involves $(n-1)^2$ operations,
and these are all parallel, so you could theoretically use $(n-1)^2$ processing elements
(processors, cores, simd lines, etc.)

\begin{exercise}
  \label{ex:gauss-n2parallel}
  Continue this reasoning.
  With $p=n^2$ processing elements each of the $(i,j)$ updates in the subblock
  can be done simultaneously.
  To be precise, how long does an arbitrary $k$ iteration take?
  Summing over all~$k$, what is the resulting $T_p, S_p,E_p$?
  How does this relate to the bounds you derived above?

  Also, with $p=n$ processing elements you could let each row or column of the subblock
  update be done in parallel. What is now the time for the $k$th outer iteration?
  What is the resulting $T_p, S_p,E_p$?,
\end{exercise}

So far we've looked at parallel execution aspects,
but without worrying about memory. 
For a memory analysis we establish the following relations
for time and memory in terms of problem size~$N$
and processor frequency~$f$:
\[ T=\frac 13 N^3/f,\qquad M=N^2. \]

Note that this formula holds for a single processor computer.

\begin{exercise}
  \label{ex:gauss-f2f}
  Suppose you buy a processor twice as fast, and you want to do a
  benchmark run that again takes time~$T$. How much memory do you need?
\end{exercise}

With $p$ parallel processing elements,
the relations become
\[ T=\frac 13 N^3/(pf),\qquad M=N^2. \]

Let's now investigate the implications of
an increase in the number of processors,
where we assume that each processor has its own memory,
containing $M_p=N^2/p$ of the total problem.
(Since keep the process itself constant, we can set
for simplicity~$f=1$.)

\begin{exercise}
  \label{ex:gauss-p2p}
  Suppose you have a cluster with $p$ processors, each with $M_p$
  memory, can run a Gaussian elimination of an $N\times N$
  matrix in time~$T$:
  \[ T=\frac 13 N^3/p,\qquad M_p=N^2/p. \]

  Now you extend the cluster to $2P$ processors, of the same clock speed,
  and you want to do a benchmark run, again taking time~$T$.
  How much memory does each node need?

  Hint: for the extended cluster:
  \[ T'=\frac 13 N'^3/p',\qquad M_p'=N'^2/p'. \]
  The question becomes to compute $M_p'$
  under the given conditions.
\end{exercise}

\Level 1 {Factorization, implementation of the dense case}

A full analysis of the scalability of
%
\emph{parallel dense LU factorization}
%
is quite
involved, so we will state without further proof that as in the matrix-vector case a
two-dimensional distribution is needed.

\begin{exercise}
  Analyze the running time, speedup, and efficiency as a function
  of~$N$, if we assume a one-dimensional distribution, and enough
  processors to store one column per processor. Show that speedup is
  limited.

  Also perform this analysis for a two-dimensional decomposition where
  each processor stores one element.
\end{exercise}

\index{matrix!storage, cyclic|(}

For this reason, an \indexterm{overdecomposition} is used, where
the matrix is divided in more blocks than there are processors,
and each processor stores several, non-contiguous, sub-matrices. 
\begin{figure}[ht]
  \includegraphics[scale=.11]{cyclic-1}
  \caption{One-dimensional cyclic distribution: assignment of four
    matrix columns to two processors, and the resulting mapping of
    storage to matrix columns.}
  \label{fig:cyclic-1}
\end{figure}
We illustrate this in figure~\ref{fig:cyclic-1} where we divide four
block columns of a matrix to two processors: each processor stores in
a contiguous block of memory two non-contiguous matrix columns.

Next, we illustrate in figure~\ref{fig:cyclic-1-mvp}
\begin{figure}[ht]
  \includegraphics[scale=.11]{cyclic-1-mvp}
  \caption{Matrix-vector multiplication with a cyclicly distributed matrix.}
  \label{fig:cyclic-1-mvp}
\end{figure}
that a matrix-vector product with such a matrix can be performed
without knowing that the processors store non-contiguous parts of the
matrix. All that is needed is that the input vector is also cyclicly
distributed.

\begin{exercise}
  Now consider a $4\times4$ matrix and a $2\times2$ processor
  grid. Distribute the matrix cyclicly both by rows and columns. Show
  how the matrix-vector product can again be performed using the
  contiguous matrix storage, as long as the input is distributed
  correctly. How is the output distributed? Show that more
  communication is needed than the reduction of the one-dimensional
  example.
\end{exercise}

Specifically,
with $P<N$ processors, and assuming for simplicity $N=cP$, we let
processor~0 store rows $0,c,2c,3c,\ldots$; processor~1 stores rows
$1,c+1,2c+1,\ldots$, et cetera. This scheme can be generalized two a
two-dimensional distribution, if $N=c_1P_1=c_2P_2$
and~$P=P_1P_2$. This is called a 2D \indexterm{cyclic
  distribution}. This scheme can be further extended by considering
block rows and columns (with a small block size), and assigning to
processor~0 the \emph{block} rows $0,c,2c,\ldots$.

\begin{exercise}
  Consider a square $n\times n$ matrix, and a square $p\times p$
  processor grid, where $p$ divides~$n$ without remainder. Consider
  the overdecomposition outlined above, and make a sketch of matrix
  element assignment for the specific case $n=6,p=2$. That is, draw an
  $n\times n$ table where location $(i,j)$ contains the processor
  number that stores the corresponding matrix element. Also make a
  table for each of the processors describing the local to global mapping,
  that is, giving the global $(i,j)$ coordinates of the elements in
  the local matrix. (You will find this task facilitated by using
  zero-based numbering.)

  Now write functions $P,Q,I,J$ of $i,j$ that describe the global to
  local mapping, that is, matrix element $a_{ij}$ is stored in
  location $(I(i,j),J(i,j))$ on processor $(P(i,j),Q(i,j))$.
\end{exercise}

\index{matrix!storage, cyclic|)}

\index{LU factorization!computation in parallel!dense|)}

\Level 1 {Factorization, sparse case}
\index{LU factorization!computation in parallel!sparse|(}

\emph{Sparse matrix LU factorization in parallel} is a notoriously
difficult topic. Any type of factorization involves a sequential
component, making it non-trivial to begin with. To see the problem
with the sparse case in particular, suppose you process
the matrix rows sequentially. The dense case has enough elements per
row to derive parallelism there, but in the sparse case that number
may be very low.

The way out of this problem is to realize that we are not interested
in the factorization as such, but rather that we can use it to solve a
linear system. Since permuting the matrix gives the same solution,
maybe itself permuted, we can explore permutations that have a higher
degree of parallelism.

The topic of matrix orderings already came up in
section~\ref{sec:fill-ordering}, motivated by fill-in reduction.
We will consider orderings with favorable parallelism properties
below: nested dissection in section~\ref{sec:dissection}, and
multi-color orderings in section~\ref{sec:redblackgreen}.

\index{LU factorization!computation in parallel!sparse|)}

